\section{Identification and sensitivity}\label{sec:identification}

This section asks what the data can and cannot say about the dark figure under \cref{asm:fw-closed,asm:fw-linkage,asm:fw-scope}, and what \cref{asm:fw-dependence} adds. Without a restriction on $\omega$, the dark figure is not identified (\cref{thm:nonid}). With $\omega$ fixed, the maximum likelihood estimator has a closed form in which $\omega$ is an explicit multiplicative sensitivity parameter (\cref{thm:closed}). This yields sharp bounds (\cref{cor:bounds}) and the direction of bias (\cref{prop:direction}), including under heterogeneous visibility (\cref{prop:id-hetero}). We then recall Chao's lower bound, which holds under heterogeneity across incidents (model $M_h$), and show that heterogeneity across lists reverses it (\cref{thm:chao,prop:id-mt}). Finally, we model referral between an NGO audit and police records (\cref{prop:referral}) and characterise the sensitivity frontier of a policy-relevant claim (\cref{thm:id-frontier}). Proofs, and two further remarks on the closed form, are in \cref{app:idproofs}.

Several of these results are classical, and we attribute them where they are stated. They are the non-estimability of the $K$-way interaction and the closed-form estimator without it, an alternating product of cell counts valid for every $K$ \citep{fienberg1972,bishop1975}; the Lincoln--Petersen estimator \citep{petersen1896,lincoln1930}; and Chao's bound \citep{chao1987}. \Citet{aleshinguendel2024} treat the absence of the highest-order interaction as an explicit identifying assumption for general $K$, separate from the model for the observed cells; their analysis overlaps with \cref{lem:id-param,thm:nonid}. Offset-based sensitivity analysis of log-linear population-size estimators was proposed by \citet{gerritse2015}, and retaining $\omega$ in the closed form simply makes such an offset explicit. Our contribution is to keep $\omega$ as an explicit sensitivity parameter and to draw out its consequences: sharp bounds with inference (\cref{cor:bounds}), the sensitivity frontier of a policy claim (\cref{thm:id-frontier}), and an interpretation of the sign of $\omega$ through heterogeneous visibility (\cref{prop:id-hetero}) and referral (\cref{prop:referral}). Along the way we note that list heterogeneity reverses Chao's inequality (\cref{prop:id-mt}). All asymptotic statements refer to the following regime.

\begin{definition}[Large-count regime]\label{def:id-regime}
Under (P), $\mu_{\h}=\tau m_{\h}$ for a fixed $\bm m\in(0,\infty)^{\Hset}$ and $\tau\to\infty$. Under (M), $\bm\pi$ is fixed and $N\to\infty$; we then set $\tau=N$ and $\bm m=\bm\pi$. In both cases $\pi_{\zero}=\mu_{\zero}/\sum_{\h\in\Hset}\mu_{\h}$ is the probability that an incident is missed by every list.
\end{definition}

\subsection{The dark figure is not identified without a restriction on \texorpdfstring{$\omega$}{omega}}\label{sec:id-nonid}

The engine of this section is a reparameterisation. Let $\Theta_\omega=\{\bm u\in\mathbb R^{2^K}:u_{\Kset}=\omega\}$, and let $\bm\mu_+(\bm u)=(\exp\sum_{S\subseteq\supp(\h)}u_S)_{\h\in\Hobs}$ be the observed-cell means implied by $\bm u$.

\begin{lemma}[Saturation given $\omega$]\label{lem:id-param}
For every $\omega\in\mathbb R$, the map $\bm u\mapsto\bm\mu_+(\bm u)$ restricted to $\Theta_\omega$ is a bijection onto $(0,\infty)^{\Hobs}$. Its inverse sends $\bm\mu_+$ to the vector $\bm u$ given by \eqref{eq:fw-mobius} with $\mu_{\zero}=e^{(-1)^K\omega}\mathcal A(\bm\mu_+)$.
\end{lemma}

In words, once $\omega$ is fixed the model for the observed cells is saturated: $2^K-1$ free parameters for $2^K-1$ cells. Each value of $\omega$ is therefore compatible with every observed-cell mean vector, and \eqref{eq:fw-key} turns the choice of $\omega$ into a multiplicative choice of the dark figure.

\begin{theorem}[Non-identification]\label{thm:nonid}
Let the log-linear model be saturated, with no restriction on $\omega$.
\begin{enumerate}[(a)]
  \item \emph{Poisson model.} Let $\bm u$ have top interaction $\omega$. For every $\omega'\in\mathbb R$ there is exactly one $\bm u(\omega')\in\Theta_{\omega'}$ with $\bm\mu_+(\bm u(\omega'))=\bm\mu_+(\bm u)$. Along this curve,
  \begin{equation}\label{eq:id-curve}
    \mu_{\zero}\bigl(\bm u(\omega')\bigr)=e^{(-1)^K(\omega'-\omega)}\,\mu_{\zero}(\bm u),
    \qquad
    u_S(\omega')=u_S+(-1)^{K-|S|}(\omega'-\omega)\quad(S\subseteq\Kset).
  \end{equation}
  Under (P), every point of the curve induces the same law of the observed table, and the observed-data likelihood \eqref{eq:fw-likP} is constant in $\omega'$. As $\omega'$ ranges over $\mathbb R$, $\mu_{\zero}$ ranges over $(0,\infty)$ and $\E N=\mu_{\mathrm{obs}}+\mu_{\zero}$ over $(\mu_{\mathrm{obs}},\infty)$. No $u$-term is identified, not even a main effect.
  \item \emph{Multinomial model.} Under (M), the conditional law of $\bm n_+$ given $n$ depends on $\bm\pi$ only through $\bm\pi^c$. For each $\omega$, the map $\bm\pi\mapsto\bm\pi^c$ is a bijection from $\{\bm\pi:u_{\Kset}=\omega\}$ onto the interior of the simplex on $\Hobs$. Moreover, $N$ is not consistently estimable: no estimator $\Nhat=\Nhat(\bm n_+)$ satisfies $\Nhat/N\to1$ in probability as $N\to\infty$ for every fixed $\bm\pi$.
\end{enumerate}
\end{theorem}

Part (a) follows from \cref{lem:id-param}: fix $\bm\mu_+$, vary $\omega$, and read off $\mu_{\zero}$ from \eqref{eq:fw-key}. For (b), the observed table informs $(N,\pi_{\zero})$ only through the count $n\sim\Bin(N,1-\pi_{\zero})$ (\cref{lem:fw-factor}). Two parameter points with the same $\bm\pi^c$ and matched binomial means but different $N$ cannot be separated, by the local limit theorem (\cref{app:id-nonid}).

\begin{remark}[The restriction on $\omega$ cannot be tested]\label{rem:id-untestable}
Every goodness-of-fit statistic, information criterion and posterior predictive check is a function of the observed table. Its law is constant along the curve in \cref{thm:nonid}(a). Model search over log-linear models (\cref{sec:models}) therefore chooses among restrictions on lower-order terms \emph{given} a restriction on $\omega$, never the restriction on $\omega$ itself. A non-hierarchical model that sets some lower-order terms to zero and leaves $\omega$ free can identify $\omega$, because the observed table then over-determines the remaining parameters. Each such restriction, however, is itself a statement about the unobservable cell, since every $u_S$ involves $\log\mu_{\zero}$ through the $T=\emptyset$ term of \eqref{eq:fw-mobius}. That is identification by functional form, not by data (compare \citealp{fienberg1972,international1995}).
\end{remark}

\subsection{A closed form with a sensitivity parameter}\label{sec:id-closed}

For $\omega=0$, the closed form below is classical for every $K$ \citep{fienberg1972,bishop1975}. We state it with $\omega$ retained.

\begin{theorem}[Closed-form estimator]\label{thm:closed}
Suppose \cref{asm:fw-closed,asm:fw-linkage,asm:fw-scope} hold and \cref{asm:fw-dependence}(a) holds with known $\omega_0=\omega$.
\begin{enumerate}[(a)]
  \item Under (P), the maximum likelihood estimator of $\mu_{\zero}$ exists if and only if $n_{\h}>0$ for every $\h\in\Hobs$. It is then
  \begin{equation}\label{eq:id-closed}
    \widehat n_{\zero}(\omega)=e^{(-1)^K\omega}\,\mathcal A(\bm n_+)
    =e^{(-1)^K\omega}\,\frac{\prod_{|\h|\ \mathrm{odd}}n_{\h}}{\prod_{|\h|\ \mathrm{even},\,\h\neq\zero}n_{\h}},
  \end{equation}
  the fitted observed cells equal the observed counts, $\widehat\mu_{\h}=n_{\h}$ for $\h\in\Hobs$, and $\Nhat(\omega)=n+\widehat n_{\zero}(\omega)$.
  \item Under (M), the conditional maximum likelihood estimator of $\bm\pi^c$ exists in the interior of the simplex if and only if $n_{\h}>0$ for every $\h\in\Hobs$, and it is then $\bm n_+/n$. Substituting it into \eqref{eq:fw-keyM} gives $\widehat\pi_{\zero}$ with $n/(1-\widehat\pi_{\zero})=\Nhat(\omega)$. Over integers $N\ge n$, the binomial factor in \eqref{eq:fw-likM} at $\pi_{\zero}=\widehat\pi_{\zero}$ attains its maximum at $\lfloor\Nhat(\omega)\rfloor$, which is the conditional maximum likelihood estimator of \citet{sanathanan1972}.
  \item For $K=2$, $\widehat n_{\zero}(\omega)=e^{\omega}n_{10}n_{01}/n_{11}$, where $e^\omega$ is the odds ratio between the two lists. At $\omega=0$, $\Nhat(0)=n_{1+}n_{2+}/n_{11}$, the Lincoln--Petersen estimator.
  \item In the regime of \cref{def:id-regime}, let $v_+=\sum_{\h\in\Hobs}\mu_{\h}^{-1}$ and $v=v_++\mu_{\zero}^{-1}=\sum_{\h\in\Hset}\mu_{\h}^{-1}$. Then:
  \begin{enumerate}[(i)]
    \item under (P), $(\log\widehat n_{\zero}(\omega)-\log\mu_{\zero})/v_+^{1/2}\Rightarrow\mathrm N(0,1)$, so that $\Var\{\log\widehat n_{\zero}(\omega)\}\approx\sum_{\h\in\Hobs}\mu_{\h}^{-1}$; under (M) the same holds with $v_+$ replaced by $v_+-N^{-1}$;
    \item under (P) and under (M), $(\log\widehat n_{\zero}(\omega)-\log\nzero)/v^{1/2}\Rightarrow\mathrm N(0,1)$. Equivalently, the prediction error $\Nhat(\omega)-N=\widehat n_{\zero}(\omega)-\nzero$ has asymptotic variance $\mu_{\zero}^2v_++\mu_{\zero}$;
    \item $\widehat v_+=\sum_{\h\in\Hobs}n_{\h}^{-1}$ and $\widehat v=\widehat v_++\widehat n_{\zero}(\omega)^{-1}$ satisfy $\widehat v_+/v_+\to1$ and $\widehat v/v\to1$ in probability.
  \end{enumerate}
\end{enumerate}
\end{theorem}

By \cref{lem:id-param}, with $\omega$ fixed the model for the observed cells is the full positive orthant, so the likelihood is maximised cell by cell at $\widehat\mu_{\h}=n_{\h}$, and \eqref{eq:fw-key} carries this maximiser to $\mu_{\zero}$. Part (d) holds because $\log\widehat n_{\zero}$ is linear in the cell logarithms with coefficients $s_{\h}=\pm1$ (\cref{app:id-closed}).

\begin{remark}[Intervals]\label{rem:id-intervals}
\Cref{thm:closed}(d) gives two Wald intervals on the log scale. For the expected dark figure under (P), $\widehat n_{\zero}\exp(\pm z_{1-\alpha/2}\widehat v_+^{1/2})$ has asymptotic coverage $1-\alpha$. For the realised total under (P) or (M), the interval
\[
  \bigl[\,n+\widehat n_{\zero}\,e^{-z_{1-\alpha/2}\widehat v^{1/2}},\;n+\widehat n_{\zero}\,e^{z_{1-\alpha/2}\widehat v^{1/2}}\,\bigr]
\]
has asymptotic coverage $1-\alpha$, and its lower end never falls below $n$. This is in the spirit of the log-transformed interval of \citet{chao1987}. For $K=2$ and $\omega=0$, a direct computation gives $\mu_{\zero}^2v_++\mu_{\zero}=\mu_{1+}\mu_{2+}\mu_{10}\mu_{01}/\mu_{11}^3$, the familiar large-sample variance of the Lincoln--Petersen estimator \citep{bishop1975}. For the illustrative \cref{tab:fw-example}, $\widehat v_+=0.104$, the $95\%$ interval for $\mu_{\zero}$ is $[382,1355]$, and the interval for $N$ is $[1141,2121]$.
\end{remark}

Empty cells, which are the rule in small domains, are handled as in \cref{rem:id-zeros}. \Cref{rem:id-unify} shows that every hierarchical log-linear estimator is the alternating product $\mathcal A$ of its \emph{fitted} cells, so these estimators differ only in how they smooth the table before applying $\mathcal A$ (both remarks are in \cref{app:id-remarks}).

\subsection{Sharp bounds and the direction of bias}\label{sec:id-bounds}

\begin{corollary}[Sharp sensitivity bounds]\label{cor:bounds}
Suppose \cref{asm:fw-closed,asm:fw-linkage,asm:fw-scope} hold and \cref{asm:fw-dependence}(b) holds with bound $\Omega\ge0$.
\begin{enumerate}[(a)]
  \item \emph{Identified set.} Under (P), the set of values of $\mu_{\zero}$ compatible with the law of the observed table and with $|\omega|\le\Omega$ is
  \[
    \bigl[\,e^{-\Omega}\mathcal A(\bm\mu_+),\;e^{\Omega}\mathcal A(\bm\mu_+)\,\bigr].
  \]
  Each point is attained by exactly one parameter value with $|\omega|\le\Omega$, so the bounds are sharp. Under (M), the same holds for the odds $\pi_{\zero}/(1-\pi_{\zero})$ of being missed by every list, given the conditional law of $\bm n_+$ given $n$, with $\mathcal A(\bm\pi^c)$ in place of $\mathcal A(\bm\mu_+)$. By \cref{thm:nonid}(b), the remaining binomial factor cannot consistently narrow this set.
  \item \emph{Estimated bounds.} If all $n_{\h}>0$, then $\{\widehat n_{\zero}(\omega):|\omega|\le\Omega\}=[e^{-\Omega}\widehat n_{\zero}(0),\,e^{\Omega}\widehat n_{\zero}(0)]$, and the bounds for $N$ are $n+e^{\pm\Omega}\widehat n_{\zero}(0)$.
  \item \emph{Inference.} In the regime of \cref{def:id-regime} under (P), for every $\omega\in[-\Omega,\Omega]$ the interval
  \begin{equation}\label{eq:id-im}
    \bigl[\,\widehat n_{\zero}(0)\,e^{-\Omega-c\,\widehat v_+^{1/2}},\;\widehat n_{\zero}(0)\,e^{\Omega+c\,\widehat v_+^{1/2}}\,\bigr],
    \qquad
    \Phi\bigl(c+2\Omega\,\widehat v_+^{-1/2}\bigr)-\Phi(-c)=1-\alpha,
  \end{equation}
  contains $\mu_{\zero}$ with probability tending to at least $1-\alpha$. Replacing $c$ by $z_{1-\alpha/2}$ gives the union of the Wald intervals over $|\omega|\le\Omega$, which is conservative.
\end{enumerate}
\end{corollary}

Part (c) is the construction of \citet{imbens2004} for a partially identified parameter \citep[see also][]{manski2003}. It applies cleanly here because the pivot $\log\widehat n_{\zero}(\omega)-\log\mu_{\zero}=\sum_{\h}s_{\h}\log(n_{\h}/\mu_{\h})$ does not involve $\omega$. The width of the identified set is $2\Omega$ on the log scale whatever $K$ is. Adding lists does not shrink the sensitivity to the top interaction; it only changes what that interaction means. For \cref{tab:fw-example} with $\Omega=\log 2$, the estimated bounds are $\widehat n_{\zero}\in[360,1440]$ and $N\in[1120,2200]$, and the interval \eqref{eq:id-im} for $\mu_{\zero}$ is $[212,2449]$. By \eqref{eq:fw-slice} with $k=3$, $\Omega=\log2$ says that the police--NGO odds ratio among incidents missed by the community security organisation is within a factor of two of the odds ratio among those it records.

\begin{proposition}[Direction of bias]\label{prop:direction}
Suppose \cref{asm:fw-closed,asm:fw-linkage,asm:fw-scope} hold and the true top interaction is $\omega$, with oriented value $\bar\omega=(-1)^K\omega$, but the analyst imposes $\omega=0$. In the regime of \cref{def:id-regime}, under (P) or (M),
\begin{equation}\label{eq:id-direction}
  \frac{\widehat n_{\zero}(0)}{\nzero}\to e^{-\bar\omega},
  \qquad
  \frac{\Nhat(0)}{N}\to1-\pi_{\zero}\bigl(1-e^{-\bar\omega}\bigr)
  \qquad\text{in probability}.
\end{equation}
Hence $\widehat n_{\zero}(0)$ asymptotically underestimates the dark figure if $\bar\omega>0$ and overestimates it if $\bar\omega<0$. For every list $k$, \eqref{eq:fw-slice} gives $\bar\omega=(-1)^{K-1}\{\psi_k(0)-\psi_k(1)\}$: imposing $\omega=0$ replaces the unobservable interaction $\psi_k(0)$ by the observable $\psi_k(1)$. In particular:
\begin{enumerate}[(a)]
  \item For $K=2$, $\bar\omega=\omega=\log\mathrm{OR}_{12}$. Lincoln--Petersen underestimates the dark figure when the lists are positively dependent ($\mathrm{OR}_{12}>1$), for example when incidents recorded by police are more likely to reach the NGO audit, and overestimates it under negative dependence.
  \item For $K=3$, $\bar\omega=-\omega=\log\mathrm{OR}_{12\mid h_3=0}-\log\mathrm{OR}_{12\mid h_3=1}$. The estimator underestimates if and only if the odds ratio between two lists is larger among incidents the third list misses than among incidents it records. Since $\omega$ is symmetric, this holds for every choice of the third list. Positive pairwise dependence alone does not determine the sign, because the $\omega=0$ model estimates all pairwise interactions freely.
  \item For general $K$, a given sign of $\omega$ has opposite consequences at odd and even $K$. We therefore report $\bar\omega$ rather than $\omega$: $e^{\bar\omega}$ is the factor by which the $\omega=0$ estimate of the dark figure must be multiplied.
\end{enumerate}
\end{proposition}

The direction in (a) is the classical warning about dependent lists \citep{international1995,wolter1986}. The point of (b) and (c) is that this intuition does not carry over: with three or more lists, ``the lists are positively dependent'' says nothing about the sign of the bias. Heterogeneous visibility is a natural source of dependence in incident data. Highly visible incidents (violent attacks, incidents with media coverage) are likely to be recorded by many lists, and inconspicuous ones by few or none. The next result shows that its effect on the sign is governed by the shape of the visibility distribution.

\begin{proposition}[Heterogeneous visibility and the sign of $\omega$]\label{prop:id-hetero}
Suppose that, under (M), each incident has a capture probability $p_i$ drawn independently from a non-degenerate distribution $F$ supported in $[\eta,1-\eta]$ for some $\eta\in(0,\tfrac12)$. Suppose also that, given $p_i$, the $K$ lists record incident $i$ independently, each with probability $p_i$; this is model $M_h$ of \citet{otis1978}. Let $X=\logit(p)$. For $t\in\mathbb R$, let $G_t$ be the law of $X$ when $p$ has density proportional to $p^t(1-p)^{K-t}$ with respect to $F$, and let $\kappa_r(t)$ be the $r$th cumulant of $G_t$. Then
\begin{equation}\label{eq:id-irwinhall}
  \omega=\E\{\kappa_K(S_K)\},\qquad S_K=U_1+\dots+U_K,\quad U_j\overset{\mathrm{iid}}{\sim}\mathrm{Uniform}(0,1),
\end{equation}
so $S_K$ has the Irwin--Hall distribution. In particular:
\begin{enumerate}[(a)]
  \item for $K=2$, $\omega=\E\{\Var_{G_{S_2}}(X)\}>0$, so Lincoln--Petersen underestimates the dark figure under every non-degenerate heterogeneity of this kind;
  \item for $K\ge3$, the sign of $\omega$ is that of an average of $K$th cumulants of the tilted laws $G_t$, and either sign can occur. By \cref{prop:direction}, the $\omega=0$ estimator underestimates if and only if $(-1)^K\E\{\kappa_K(S_K)\}>0$.
\end{enumerate}
\end{proposition}

Part (a) recovers the classical result that heterogeneity biases Lincoln--Petersen downward \citep{international1995,chao2001}. We show that the general-$K$ case is governed by higher cumulants: skewness of the tilted logit-visibility distribution at $K=3$, and kurtosis at $K=4$.

\begin{example}[Visibility heterogeneity can reverse the bias]\label{ex:id-hetero}
Suppose $95\%$ of incidents have $p=0.1$ and $5\%$ are highly visible with $p=0.9$. Computing $\omega$ exactly from $\pi_{\h}=\int p^{|\h|}(1-p)^{K-|\h|}dF$ gives $\omega=1.56$ for $K=2$, $2.53$ for $K=3$ and $-0.49$ for $K=4$. By \eqref{eq:id-direction}, the limit of $\widehat n_{\zero}(0)/\nzero$ is $e^{-\bar\omega}$: $0.21$ for $K=2$, $12.5$ for $K=3$ and $1.63$ for $K=4$. Lincoln--Petersen recovers a fifth of the dark figure. For $K=3$, the $\omega=0$ estimator, which is the model with all two-way interactions, overstates the dark figure $12.5$-fold and $N$ about $9$-fold. The reason is that the few highly visible incidents pile up in $n_{111}$, which enters $\mathcal A$ in the numerator. Simpler three-list log-linear models fitted to the same population understate the dark figure instead: the limit of $\widehat n_{\zero}/\nzero$ is $0.24$ under independence, $0.30$ with one two-way interaction and $0.69$ with two. The magnitude is also fragile. Over two-class mixtures with low visibility $p\in\{0.05,0.1,0.2\}$, high visibility $p\in\{0.7,0.8,0.9\}$ and a highly visible share in $\{0.02,0.05,0.1,0.2\}$, the all-two-way limit ranges from $0.37$ to $25$. These are exact population limits for stylised models, not estimates. The lesson is not that three lists are worse than two. It is that ``heterogeneity'' fixes neither the sign nor the size of the bias, which is why \cref{sec:models} models heterogeneity directly through latent classes.
\end{example}

\subsection{Chao's bound and its reversal}\label{sec:id-chao}

Chao's estimator is often described as robust to heterogeneity. It is a lower bound under one kind of heterogeneity, across incidents, and it can fail in the other direction under heterogeneity across lists.

\begin{theorem}[Chao's lower bound under $M_h$; \citealp{chao1987}]\label{thm:chao}
Suppose that, conditionally on $p_1,\dots,p_N\in[0,1]$, the indicators $Y_{ik}$ are independent with $\Pr(Y_{ik}=1)=p_i$ for every list $k$ (model $M_h$). Write $\E$ for expectation conditional on $\bm p$, and suppose $\E f_2>0$. Then
\begin{equation}\label{eq:id-chao}
  \E f_0\;\ge\;\frac{K-1}{K}\,\frac{(\E f_1)^2}{2\,\E f_2},
  \qquad\text{equivalently}\qquad
  N\;\ge\;\E n+\frac{K-1}{K}\,\frac{(\E f_1)^2}{2\,\E f_2}.
\end{equation}
If every $p_i\in(0,1)$, equality holds if and only if $p_1=\dots=p_N$. If, in addition, the $p_i$ are i.i.d.\ from a distribution $F$ on $(0,1)$, then the estimator
\[
  \Nhat_{\mathrm C}=n+\frac{K-1}{K}\,\frac{f_1^2}{2f_2}
\]
satisfies $\Nhat_{\mathrm C}/N\to\lambda_{\mathrm C}\le1$ almost surely as $N\to\infty$, with $\lambda_{\mathrm C}=1$ if and only if $F$ is a point mass.
\end{theorem}

The bound uses only the singletons $f_1$ and doubletons $f_2$, the incidents that are hardest to record, and its proof is a single application of Cauchy--Schwarz to $\sum_i p_i(1-p_i)^{K-1}$. It holds whatever the distribution of the incident-level $p_i$, but only under $M_h$, which requires the lists to be exchangeable for each incident. That is implausible when a police list and a community portal differ greatly in reach, and heterogeneity \emph{across lists} pushes Chao's estimator in the opposite direction.

\begin{proposition}[List heterogeneity reverses Chao's inequality]\label{prop:id-mt}
Suppose instead that the $Y_{ik}$ are independent with $\Pr(Y_{ik}=1)=a_k\in(0,1)$ for every incident (model $M_t$). Then
\[
  \E f_0\;\le\;\frac{K-1}{K}\,\frac{(\E f_1)^2}{2\,\E f_2},
\]
with equality if and only if $a_1=\dots=a_K$. Consequently, $\Nhat_{\mathrm C}/N$ converges almost surely to a limit $\ge1$: under $M_t$, Chao's estimator is an asymptotic \emph{upper} bound.
\end{proposition}

Under $M_t$ the independence model of \citet{darroch1958} holds, so $\omega=0$ and \cref{thm:closed} is consistent. When incidents and lists are both heterogeneous, the two effects compete and $\Nhat_{\mathrm C}$ can err in either direction. Chao's estimator is therefore a lower bound under $M_h$, not under heterogeneity in general, and we treat it as a diagnostic. The only floor that needs no assumption is the observed count $n$. For \cref{tab:fw-example}, $\Nhat_{\mathrm C}=760+\tfrac23\cdot540^2/380=1272$, below the $\omega=0$ estimate of $1480$. Since the three illustrative lists differ in size, \cref{prop:id-mt} warns against reading $1272$ as a hard floor.

\subsection{Referral dependence}\label{sec:id-referral}

An NGO audit may ingest incidents that were first reported to police, for example through public police statements, data-sharing agreements, or victims who report to both and say so. This creates a specific and directional form of dependence, which we model explicitly.

\begin{proposition}[Referral dependence]\label{prop:referral}
Let $K=2$ (police $k=1$, NGO audit $k=2$), and let incidents be i.i.d.\ as follows. Police record an incident with probability $p_1\in(0,1)$, so $Y_1\sim\Bern(p_1)$. The NGO captures it directly (through a victim report, media monitoring or outreach) with probability $p_2\in(0,1)$, so $D\sim\Bern(p_2)$. If police record it, the NGO also ingests the police record with probability $q\in[0,1)$, so $R\sim\Bern(q)$. The variables $Y_1$, $D$ and $R$ are independent, and $Y_2=\max(D,RY_1)$.
\begin{enumerate}[(a)]
  \item \emph{Induced dependence.} The cell probabilities are $\pi_{11}=p_1\{p_2+q(1-p_2)\}$, $\pi_{10}=p_1(1-p_2)(1-q)$, $\pi_{01}=(1-p_1)p_2$ and $\pi_{00}=(1-p_1)(1-p_2)$. Hence
  \begin{equation}\label{eq:id-referral-or}
    \mathrm{OR}_{12}=\frac{\pi_{11}\pi_{00}}{\pi_{10}\pi_{01}}=1+\frac{q}{p_2(1-q)},
    \qquad
    u_{12}=\omega=\log\Bigl\{1+\frac{q}{p_2(1-q)}\Bigr\}\ge0,
  \end{equation}
  with equality if and only if $q=0$. The odds ratio does not depend on $p_1$. It increases in $q$ and decreases in $p_2$.
  \item \emph{Bias of Lincoln--Petersen.} As $N\to\infty$, almost surely,
  \begin{equation}\label{eq:id-referral-bias}
    \frac{\Nhat_{\mathrm{LP}}}{N}\to\frac{p_2+p_1q(1-p_2)}{p_2+q(1-p_2)}=1-\frac{q(1-p_1)(1-p_2)}{p_2+q(1-p_2)},
    \qquad
    \frac{\widehat n_{\zero}^{\mathrm{LP}}}{\nzero}\to\frac{1}{\mathrm{OR}_{12}}=\frac{p_2(1-q)}{p_2+q(1-p_2)} .
  \end{equation}
  \item \emph{Full provenance.} Suppose each NGO record shows whether the NGO captured the incident directly ($D=1$), whether or not it also ingested a police record. Recode $Y_2^\circ=D$. This drops referral-only captures ($Y_1=1$, $R=1$, $D=0$) from list $2$ and moves them from cell $11$ to cell $10$. Then $Y_1$ and $Y_2^\circ$ are independent, the recoded table has $\omega=0$, and $\Nhat^\circ=n_{1+}n^\circ_{2+}/n^\circ_{11}$ satisfies $\Nhat^\circ/N\to1$ almost surely. More generally, for any $K$: if the direct-capture histories $(Y_1,D,Y_3,\dots,Y_K)$ satisfy \cref{asm:fw-dependence}(a), so does the recoded table.
  \item \emph{Referral flags only.} Suppose only the flag $Z=RY_1$ (``this NGO record was ingested from police'') is observed, and direct capture of flagged incidents is not. Let $n_Z$ be the number of flagged incidents, $m_{11}$ the number of unflagged incidents recorded by both lists, and $n_{01}$ the number recorded by the NGO only. Remove the flagged incidents, apply Lincoln--Petersen to the rest, and add them back:
  \[
    \Nhat^{\ddagger}=n_Z+\frac{(n_{1+}-n_Z)(m_{11}+n_{01})}{m_{11}} .
  \]
  The reduced table has odds ratio one, and $\Nhat^{\ddagger}/N\to1$ almost surely.
  \item \emph{A third, independent list.} Let $K\ge3$. Suppose $(Y_1,Y_2)$ is generated as above and $(Y_3,\dots,Y_K)$ is independent of $(Y_1,Y_2)$. Then $\omega=0$; indeed $u_S=0$ whenever $S$ meets both $\{1,2\}$ and $\{3,\dots,K\}$. Hence $\widehat n_{\zero}(0)$ is consistent: referral between two lists does not bias the $K$-list estimator.
\end{enumerate}
\end{proposition}

The dark-figure bias in \eqref{eq:id-referral-bias} depends only on $(p_2,q)$. If the NGO captures $20\%$ of incidents directly and ingests $30\%$ of police records, the odds ratio is $1+0.3/(0.2\cdot0.7)\approx3.14$. Lincoln--Petersen then recovers only about $32\%$ of the dark figure, whatever the police recording rate. Parts (c) and (d) turn a metadata field into an identification strategy, which is a reason to request provenance information from every list holder (\cref{sec:plan}). Part (e) shows that a third list unconnected to the referral channel neutralises it in the model without a $K$-way term.

\begin{remark}[Calibrating $\Omega$ from referral]\label{rem:id-calibrate}
For $K=2$, referral produces only $\omega\ge0$. Suppose the analyst is willing to bound the ingestion rate, $q\le\bar q$, and the direct coverage, $p_2\ge\underline p_2$. Then \eqref{eq:id-referral-or} gives the one-sided restriction $0\le\omega\le\log\{1+\bar q/(\underline p_2(1-\bar q))\}$. For example, $\bar q=\underline p_2=0.2$ gives $e^{\Omega}=2.25$. The bounds of \cref{cor:bounds} then become $[\widehat n_{\zero}(0),\,e^{\Omega}\widehat n_{\zero}(0)]$. This translates the abstract $\omega$ into two quantities that list holders can speak to.

For $K\ge3$ the one-sided conclusion fails, because \eqref{eq:id-referral-or} is a two-list result and referral involving several lists can give $\bar\omega$ of either sign (\cref{ex:id-referral3} in \cref{app:id-remarks}). With three or more lists the analyst should therefore calibrate $\bar\omega$, which governs the direction of bias (\cref{prop:direction}), from an explicit model of the referral channels, not assume a sign for $\omega$.
\end{remark}

\subsection{The sensitivity frontier of a policy claim}\label{sec:id-frontier}

Policy discussion of hate crime can turn on qualitative claims, such as ``most incidents never reach the police'', rather than on point estimates. We ask how large a departure from $\omega=0$ such a claim can survive. For $\kappa>0$, consider the claim
\[
  \mathsf C_\kappa:\quad\mu_{\zero}>\kappa\,\mu_{1+},
\]
where $\mu_{1+}=\E n_{1+}$. In words, the expected number of incidents recorded by no list exceeds $\kappa$ times the expected police count. For $\kappa=1$, the dark figure exceeds the police count.

\begin{theorem}[Sensitivity frontier]\label{thm:id-frontier}
Define the population frontier and its estimate
\[
  \Omega^*_\kappa=\log\frac{\mathcal A(\bm\mu_+)}{\kappa\,\mu_{1+}},
  \qquad
  \widehat\Omega_{\kappa,\alpha}=\log\frac{\widehat n_{\zero}(0)}{\kappa\,n_{1+}}-z_{1-\alpha}\,\widehat s_1,
  \qquad
  \widehat s_1^{\,2}=\sum_{\h\in\Hobs}\frac1{n_{\h}}+\frac1{n_{1+}} .
\]
\begin{enumerate}[(a)]
  \item Under (P), \cref{asm:fw-closed,asm:fw-linkage,asm:fw-scope} and \cref{asm:fw-dependence}(b) with bound $\Omega$: $\mathsf C_\kappa$ holds for every parameter value that is observationally equivalent to the truth and has $|\omega|\le\Omega$ if and only if $\Omega<\Omega^*_\kappa$.
  \item In the regime of \cref{def:id-regime}, under (P) or (M), for every $\Omega\ge0$ and every parameter value with $|\omega|\le\Omega$ at which $\mathsf C_\kappa$ fails ($\mu_{\zero}\le\kappa\mu_{1+}$),
  \[
    \limsup\ \Pr\bigl(\widehat\Omega_{\kappa,\alpha}>\Omega\bigr)\le\alpha,
  \]
  with equality when $\mu_{\zero}=\kappa\mu_{1+}$ and $\bar\omega=-\Omega$.
  \item If $\Omega<\Omega^*_\kappa$, then $\Pr(\widehat\Omega_{\kappa,\alpha}>\Omega)\to1$.
\end{enumerate}
\end{theorem}

Declaring ``$\mathsf C_\kappa$ is robust to $|\omega|\le\Omega$'' whenever $\Omega<\widehat\Omega_{\kappa,\alpha}$ is therefore an asymptotically level-$\alpha$ test of the negation of $\mathsf C_\kappa$. It is valid whatever the unidentified $\omega$ in $[-\Omega,\Omega]$ is, and it is consistent. The statistic $\widehat\Omega_{\kappa,\alpha}$ is the largest such $\Omega$; it is a breakdown point in the sense of \citet{masten2020}, and the analogue for multiple-systems estimation of the E-value \citep{vanderweele2017} and of Rosenbaum's sensitivity parameter \citep{rosenbaum2002}. The key fact is again that $\log\{\widehat n_{\zero}(\omega)/n_{1+}\}-\log(\mu_{\zero}/\mu_{1+})$ is free of $\omega$. The list $k=1$ plays no special role, and any list may replace it. \citet{gerritse2015} observe that a population-size estimate is less robust when implied coverage is low. In our parameterisation this is exact: $\partial\log\Nhat(\omega)/\partial\bar\omega=\widehat n_{\zero}(\omega)/\Nhat(\omega)$, the implied share of incidents missed by every list. Claims about the dark figure itself, such as $\mathsf C_\kappa$, behave the other way. The larger the implied dark figure relative to the police count, the larger $\Omega^*_\kappa$.

\begin{example}[Illustrative numbers, continued]\label{ex:id-frontier}
For \cref{tab:fw-example}, the claim ``the dark figure exceeds the police count'' ($\kappa=1$) has plug-in frontier $\log(720/400)=0.59$, which estimates $\Omega^*_1$. Ignoring sampling error, $\omega$ would have to move the dark figure by a factor of $1.8$ to overturn it. Once sampling error is included ($\widehat s_1=0.33$, $\alpha=0.05$), $\widehat\Omega_{1,0.05}=0.05$, so the claim is barely robust at all. The weaker claim with $\kappa=\tfrac12$ has $\widehat\Omega_{1/2,0.05}=0.74$ and survives any departure up to a factor of $e^{0.74}\approx2.1$. With counts of this size, only the weaker claim could responsibly be reported.
\end{example}
